% 原幂函数第18--20、23、25--27页；也用于指数函数原稿第2--5页复习。
\begin{frame}[t]{例2：立方根函数的图象}
 函数 $y=x^{1/3}=\sqrt[3]x$ 的图象是（\quad）。
 \par\rootchoices
 \pause
 \key{过 $(0,0)$、$(1,1)$，关于原点对称。}\\
 当 $x>1$ 时，$x^{1/3}<x$；当 $0<x<1$ 时，$x^{1/3}>x$。\\
 因而应选 \ans{B}。
\end{frame}

\begin{frame}[t]{例2：由图象判断指数大小}
 \begin{columns}[T,onlytextwidth]
  \begin{column}{.53\textwidth}\centering\powerorder\end{column}
  \begin{column}{.44\textwidth}
   根据图象中 $x>1$ 的部分，比较 $m,n,p,q$ 的大小。
   \pause
   \[n<0<q<m<1<p.\]
   因而
   \[\ans{n<q<m<p}.\]
   \reason{在 $x>1$ 时，指数越大，图象越高。}
  \end{column}
 \end{columns}
\end{frame}

\begin{frame}[t]{例3：待定指数求函数值}
 若 $f(x)$ 是幂函数，且 $f(8)=16$，求 $f(9),f(64)$。
 \pause
 设 $f(x)=x^\alpha$。由 $8^\alpha=16$ 得
 \[2^{3\alpha}=2^4\quad\Longrightarrow\quad\alpha=\frac43.\]
 因此 $f(x)=x^{4/3}$，
 \[\ans{f(9)=9^{4/3}=9\sqrt[3]9},\qquad
   \ans{f(64)=(2^6)^{4/3}=256}.\]
 \reason{函数类型已知时，先设标准形式，再求参数。}
\end{frame}

\begin{frame}[t]{例4：同指数幂的大小}
 比较各组中两个数的大小：
 \[(1)\ 1.5^{1.4},\ 1.6^{1.4};\qquad
   (2)\ 1.5^{0.4},\ 1.6^{0.4}.\]
 \pause
 \begin{itemize}
  \item $y=x^{1.4}$ 在 $[0,+\infty)$ 上递增，故
   \[\ans{1.5^{1.4}<1.6^{1.4}}.\]
  \item $y=x^{0.4}$ 在 $[0,+\infty)$ 上递增，故
   \[\ans{1.5^{0.4}<1.6^{0.4}}.\]
 \end{itemize}
\end{frame}

\begin{frame}[t]{例4：负指数与比较方法}
 \[(3)\ 1.5^{-1.5},\ 1.6^{-1.5}.\]
 \pause
 $y=x^{-1.5}$ 在 $(0,+\infty)$ 上递减，且 $1.5<1.6$，故
 \[\ans{1.5^{-1.5}>1.6^{-1.5}}.\]
 \gap
 \key{用单调性直接比较时，两个自变量应处于同一单调区间。}
 \gap
 跨区间时应另作分析，例如：$-1<1$，但 $(-1)^{-1}<1^{-1}$。\\
 $y=1/x$ 虽在两侧分别递减，却不在整个定义域上递减。
\end{frame}
